\documentclass[10pt,a4paper]{amsart}
\usepackage[margin=19mm]{geometry}
\usepackage{amsmath,amssymb,mathtools}
\usepackage[scr=rsfs,cal=euler]{mathalfa}
\usepackage{xfrac}
\usepackage[unicode,hidelinks,bookmarks=false]{hyperref}
\newcommand{\ie}{i.e.\ }
\newcommand{\cf}{cf.\ }
\newcommand{\resp}{respectively\ }
\newcommand{\Subsection}{\subsection}
\providecommand{\nobreakdash}{\nobreak\hbox{-}}
\setlength{\emergencystretch}{2em}
\multlinegap0pt
\usepackage{amssymb}
\usepackage{enumerate}
\newtheorem{corollary}{Corollary}
\newtheorem{theorem}{Theorem}
\title{The dual of the Hardy space associated to the Dunkl-Schr\"odinger operator with reverse H\"older class potential}
\author{P. Athulya, S.K. Verma}
\date{Source: arXiv:2605.14456v1 (CC BY 4.0)}
\begin{document}
\maketitle

\noindent\emph{Source and licence.} The dual of the Hardy space associated to the Dunkl-Schr\"odinger operator with reverse H\"older class potential. Version arXiv:2605.14456v1 (CC BY 4.0), distributed by arXiv under the Creative Commons Attribution 4.0 International licence, http://creativecommons.org/licenses/by/4.0/, as recorded in the arXiv licence metadata for this exact version and shown on the abstract page https://arxiv.org/abs/2605.14456v1. Full licence notice: \texttt{licenses/arxiv-2605.14456-CC-BY-4.0.txt}.\par\smallskip

\noindent\textit{Editorial scope.} Counted unit: the article's theorem theorem (source0.tex lines 841--843). The article's complete original proof of it is delivered with it (source0.tex lines 845--925, one \texttt{proof} environment of 81 lines). No mathematics is written, repaired or re-derived: the statement and the proof are the article's own lines, in the article's own notation, and no retained line is substituted or re-flowed.  Retained uncounted as the source-local context the proof needs: lines 303--315; lines 423--426; lines 455--459.  Nothing was omitted from any retained span, so the package carries no omission record.  No retained line contains a citation, so the wrapper carries no bibliography and no bibliography entry.  No retained line contains a figure environment, an image-inclusion command or drawing code, so no image is delivered and the package declares no asset dependency.  Editorial formatting only, in addition: an AMS \texttt{amsart} preamble that restates the article's own declarations and macros used by the retained lines, namely the environments \texttt{corollary}, \texttt{theorem} on separate counters, together with the standard packages those lines need.  The retained environments print in this document as Corollary 1, Theorem 1 in order of appearance, whereas the article prints the same items with the numbering of its own full document.  The complete native source of the article as distributed in the arXiv e-print for this exact version is bundled unchanged as \texttt{sources/200588/source0.tex}.
\begin{corollary} \label{Corollary J-N}
Let $f \in BMO(\mathcal{L}_k)$. Then there exist positive constants $C$ and $C^*$ such that, for every $p \in [1,\infty)$, we have
\begin{align} \label{Corollary 5.4_1}
\left(\frac{1}{\omega_k(B)} \int_B |f(x)-f_B|^p w_k(x)\,dx \right)^{\frac{1}{p}}
& \leq C \|f\|_{BMO(\mathcal{L}_k)}, \quad \text{for all balls } B,
\end{align}
and
\begin{align} \label{corollary 5.4_2}
\left(\frac{1}{\omega_k(B)} \int_B |f(x)|^p w_k(x)\,dx \right)^{\frac{1}{p}}
& \leq C^* \|f\|_{BMO(\mathcal{L}_k)}, \quad \text{for all balls } B = B(x_0, r),
\end{align}
where $x_0 \in \mathbb{R}^n$ and $r \geq \rho_k(x_0)$. Moreover, the constant $C$ depends only on the doubling constant $C_{w_k}$, $p$, and $\rho_k$.
\end{corollary}

\begin{align} 
    k_t(x,y) \leq C {\left( 1+ \mathcal{E}_k(x,y)\mathcal{}   \right) ^{-N} } \exp\left(-\frac{d_\mathcal{O}(x,y)^2}{ct} \right) \left(\omega_k(B(x,\sqrt{t}))\right)^{-1}
    , \label{assumption-1} 
\end{align} where

\begin{theorem} \label{atomic decomposition} Assume that $V$ satisfies \eqref{assumption-1}. Then there exists a constant $C>0$ such that, for every $f\in L_k^1(\mathbb{R}^n)$, we have
\begin{align*}
C^{-1}\|f\|_{H_{\Tilde{\mathcal{L}}_k}^{1,\text{atm}}} \leq \|f\|_{H_{\Tilde{\mathcal{L}}_k}^1} \leq C\|f\|_{H_{\Tilde{\mathcal{L}}_k}^{1,\text{atm}}}. 
\end{align*}
\end{theorem}

\begin{theorem}
    There exists a linear isomorphism between the Banach spaces $\left( H_{\Tilde{\mathcal{L}}_k}^1\right)^*$ and $BMO(\mathcal{L}_k)$.
\end{theorem}

\begin{proof}
Let $f \in BMO(\mathcal{L}_k)$, we define $\Phi_f$ as 
\begin{align*}
    \Phi_f(h) = \int_{\mathbb{R}^n} f(x)h(x)w_k(x)dx \quad \text{for } h\in H_{\Tilde{\mathcal{L}}_k}^1.
\end{align*}
By the atomic decomposition of $H_{\Tilde{\mathcal{L}}_k}^1$, it suffices to check the boundedness of  $\Phi_f$ on $H_{\Tilde{\mathcal{L}}_k}^{1} \text{-atoms}$. Let $a$ be an atom with $\operatorname{supp}(a)\subseteq B_0=B(x_0,r_0)$. \\

\noindent  If $r_0 < \rho_k(x_0)$, then using the cancellation property of $a$ we obtain
\begin{align*}
|\Phi_f(a)|
&= \left| \int_{\mathbb{R}^n} f(x)a(x)w_k(x)\,dx \right|
 = \left| \int_{B_0}\big(f(x)-f_{B_0}\big)a(x)w_k(x)\,dx \right| \\
&\le \frac{1}{\omega_k(B_0)} \int_{B_0} |f(x)-f_{B_0}|\, w_k(x)\,dx
\le \|f\|_{BMO(\mathcal{L}_k)} .
\end{align*}
On the other hand, if $r_0 \geq \rho_k(x_0)$, we use the size condition of $a$ and obtain
\begin{align*}
|\Phi_f(a)|
= \left| \int_{\mathbb{R}^n} f(x)a(x)w_k(x)\,dx \right|
\le \frac{1}{\omega_k(B_0)} \int_{B_0} |f(x)|\, w_k(x)\,dx
\le \|f\|_{BMO(\mathcal{L}_k)} .
\end{align*}
Since finite linear combinations of atoms are dense in $H_{\Tilde{\mathcal{L}}_k}^1$, it follows that $\Phi_f$ extends to a continuous linear functional on $H_{\Tilde{\mathcal{L}}_k}^1$, with
\[
\|\Phi_f\| \leq \|f\|_{BMO(\mathcal{L}_k)}.
\]
We now consider $\Phi \in \left(H_{\Tilde{\mathcal{L}}_k}^1\right)^*$. Let $B_0 = B(x_0,r_0)$ with $r_0 \ge \rho_k(x_0)$, and let $g \in L_k^2(B_0)$ be nonzero. Define
\[
a = \frac{g}{\|g\|_{L_k^2(B_0)}\, \omega_k(B_0)^{1/2}} .
\]
Then $a$ is a $(1,2)$-atom. Since $\Phi$ is continuous, we have
\[
|\Phi(a)| \le \|a\|_{H_{\Tilde{\mathcal{L}}_k}^1}\,\|\Phi\|.
\]
Hence, for every $g \in L_k^2(B_0)$,
\[
|\Phi(g)| \le C \|\Phi\| \|g\|_{L_k^2(B_0)} \omega_k(B_0)^{1/2},
\]
where $C$ is the constant appearing in Theorem~\ref{atomic decomposition}. It follows that $\Phi$ defines a bounded linear functional on $L_k^2(B_0)$ with operator norm bounded by $C\|\Phi\|\omega_k(B_0)^{1/2}$. By the Riesz representation theorem, there exists $f_\Phi^{B_0} \in L_k^2(B_0)$ such that
\[
\Phi(g) = \int_{B_0} g(x) f_\Phi^{B_0}(x) w_k(x)\,dx,
\]
and
\[
\|f_\Phi^{B_0}\|_{L_k^2(B_0)} \le C \|\Phi\| \omega_k(B_0)^{1/2}.
\]
\noindent If $B(y,r_1) \subset B(y,r_2)$ with $r_1 \ge \rho_k(y)$, then for every $g \in L_k^2(B(y,r_1))$,
\[
\Phi(g)
= \int_{B(y,r_1)} g(x) f_\Phi^{B(y,r_1)}(x) w_k(x)\,dx
= \int_{B(y,r_1)} g(x) f_\Phi^{B(y,r_2)}(x) w_k(x)\,dx.
\]
Hence $f_\Phi^{B(y,r_1)} = f_\Phi^{B(y,r_2)}$ a.e. on $B(y,r_1)$. Consequently, there exists a unique function $f \in L_{k,\mathrm{loc}}^2(\mathbb{R}^n)$ such that
\[
\Phi(g) = \int_B g(x) f(x) w_k(x)\,dx
\]
for every $g \in L_k^2(\mathbb{R}^n)$ with $\operatorname{supp}(g)\subset B = B(y,r)$ and $r \ge \rho_k(y)$. In particular, this holds for all $(1,\infty)$-atoms and their finite linear combinations. Hence, $\Phi$ admits an integral representation on a dense subset and extends to the whole space $H_{\Tilde{\mathcal{L}}_k}^1$. It remains to show that $f \in BMO(\mathcal{L}_k)$. Let $B = B(y,r)$ with $r \geq \rho_k(y)$. Then
\[
\frac{1}{\omega_k(B)} \int_B |f(x)| w_k(x)\,dx
\le \frac{\|f\|_{L_k^2(B)}}{\omega_k(B)^{1/2}}
\le C \|\Phi\|.
\]
For balls $B = B(y,s)$ with $s < \rho_k(y)$, take a nonzero $g \in L_k^2(B)$ satisfying
\[
\int_B g(x) w_k(x)\,dx = 0.
\]
In this case, the representative $f$ is determined by taking the modulo with constants, that is, as an element of $L_k^2(B)/\mathcal C$. Moreover,
\[
\inf_{c \in \mathcal C} \|f-c\|_{L_k^2(B)}
\le C \|\Phi\| \omega_k(B)^{1/2}.
\]
By Corollary~\ref{Corollary J-N},
\[
\frac{1}{\omega_k(B)} \int_B |f-f_B| w_k(x)\,dx
\le \left(
\frac{1}{\omega_k(B)} \int_B |f-f_B|^2 w_k(x)\,dx
\right)^{1/2}
\le C \|\Phi\|,
\]
for every ball whose radius is smaller than the critical radius. Therefore $f \in BMO(\mathcal{L}_k)$.
\end{proof}

\end{document}
